\cotefile{156-6}
\scanpage{58}\mspage{57}
on s'attend pas \inserted{\unsure{pour autant}} que $\unsure{\mathcal F|F}$ soit \unsure{une|un} sous-atelier
spacieux, au sens \deleted{le plus} fort où nous
avons introduit ce terme, i.e. stable par
\unsure{Sup|sup|soup} \uline{quelconques}. \deleted{\unsure{A|D}}\unsure{Ainsi|Aussi}, \deleted{\unsure{on}} on s'attend
pas en général que $\deleted{\unsure{\{\{x\}\}|\{x\}}}\ \{F_x=\{\{x\}\}\}_{x\in L}$
soit dans $\mathcal F$ ! La condition raisonnable
dès lors, c'est que $\mathcal F$ soit \uline{finiment}
spacieux. Vu la condition Atens${}^{\unsure{\text{é}|\prime}}$ 3,
% slanted
\margin{On dit alors que $\mathcal F$ est un atelier \unsure{ensembliste} finiment spacieux}
cela signifie simplement \deleted{\unsure{ceci}} que
\deleted{\uline{\unsure{Atensfig 3|Atens t.f. 3}}} \inserted{la relation} $\underset{M}{\mathrel{\tikz[baseline=-.45ex,x=1em,y=1em]{\draw[line width=.055em] (.52,.27)--(.05,.055)--(.48,-.11); \draw[line width=.055em] (.14,.21)--(.46,.04)--(0,-.17);}}}$ est la relation
$\underset{\mathrm{ens}}{\mathrel{\tikz[baseline=-.45ex,x=1em,y=1em]{\draw[line width=.055em] (.52,.27)--(.05,.055)--(.48,-.11); \draw[line width=.055em] (.14,.21)--(.46,.04)--(0,-.17);}}}$ \unsure{ensembliste}. Donc

\uline{Corollaire du Scholie} Les ateliers (\unsure{purs|finis})
ensemblistes finiment spacieux,
à figures de type fini, correspondent
aux $(L,M)$, $L$ un ens et $M$
un ens de figures élémentaires dans $L$,
satisfaisant \uline{\unsure{Atens t.f.|Atensfig.}} 1 et \uline{\unsure{Atens t.f|Atensfig}} 2
(i.e les $\{\{x\}\}$ sont dans $M$, et \deleted{\unsure{s|e}}si $X$ est
dans $M$ \deleted{\unsure{toute|toutes}} multistrate de $X$ l'est),
et tel que $\forall \unsure{x|X}\in M$ \deleted{\unsure{les}} \unsure{sa|son} figure
de $\unsure{\alpha|L}$ soit de t.f. Alors $\mathcal F$ est l'ens des
